\ifthenelse{\boolean{acl}}{\section{(\blt-)DP-FTRL for Private Learning}}{\section{(\blt-)DP-FTRL for Private Learning in Practice}} \label{sec:background}

\ifthenelse{\boolean{acl}}{\subsection{Background}}{}

\ifthenelse{\boolean{acl}}{
We use ($\epsilon$, $\delta$)-DP~\citep{dwork2006calibrating,dwork2014algorithmic} and $\rho$-zCDP (zero-Concentrated DP)~\citep{bun2016concentrated} to quantify the privacy protection: smaller $\epsilon$ ($\rho$) correspond to stronger DP guarantees. A formal definition and more discussion are in \cref{subsec:dp_def}.
}{
\subsection{DP Formulation} \label{subsec:dp_def}
We present the definition of ($\epsilon$, $\delta$)-DP~\citep{dwork2006calibrating,dwork2014algorithmic} to quantify the privacy protection. 
\begin{definition}[$(\epsilon,\delta)$-Differential Privacy]
  A randomized algorithm $\mathcal{M}$ satisfies  ($\epsilon$, $\delta$)-DP for $\mathbb{D}$ if for any two neighboring datasets $\mathbb{D}$, $\mathbb{D}'$ and for all $\mathcal{S}\subset \text{Range}(\mathcal{M})$: 
\[
\textstyle{\Pr[\mathcal{M}(\mathbb{D}) \in \mathcal{S}] \leq e^{\epsilon}\Pr[\mathcal{M}(\mathbb{D}') \in \mathcal{S}] +\delta}
.\]
\end{definition}
Smaller ($\epsilon$, $\delta$) values suggest stronger DP guarantees, and we often measure $\epsilon$ at a fixed small $\delta=10^{-10}$. DP-FTRL also uses an alternative definition, $\rho$-zCDP (zero-Concentrated DP)~\citep{bun2016concentrated} designed for Gaussian mechanism, and smaller $\rho$ suggests stronger DP guarantees. We use PLD (privacy Loss Distributions) accounting~\citep{doroshenko2022connect,pldlib} to convert $\rho$-zCDP to ($\epsilon$, $\delta$) DP. When applying DP in FL, the neighboring datasets $\mathbb{D}$, $\mathbb{D}'$ are defined by zeroing out the contribution of all data on a user device. More discussions on neighboring dataset for the streaming setting in learning, and connection to DP guarantees are provided in \cref{sec:sensitivity}.
\begin{algorithm}[thb]
\caption{\small FedAvg~\citep{mcmahan18learning} with \colorbox{blue!25}{DP-FTRL~\citep{kairouz21practical}} for \colorbox{blue!25}{DP} FL}
\label{algo:dpfl}
\begin{algorithmic}
\INPUT: clients per round $m$, learning rate on client $\eta_{c}$ and on server $\eta_{s}$, momentum $\beta=0.9$,  total number of rounds $T$, \colorbox{blue!25}{clip norm $\zeta$, clip norm noise multiplier $\sigma$,} 
\begin{multicols}{2}
\STATE Initialize model $y^{-1}$ with pretraining
\STATE Initialize server optimizer state $\optstate$
\STATE \colorbox{blue!25}{Initialize correlated noise state $\noisestate$ with $\sigma \zeta$}
\FOR{each round $t=0, 1, 2, \ldots, n-1$}
\STATE $\mathcal{Q}^t \leftarrow$ (at least $m$ users that did not participate in the previous $b$ rounds)
\FOR{each user $i \in \mathcal{Q}^t$ \textbf{in parallel}}
\STATE $\Delta^{t}_i \leftarrow$ ClientUpdate$(i, y^{t-1})$
\ENDFOR
\vspace{0.1cm}
\STATE \colorbox{blue!25}{$\tilde{\Delta}^t, \noisestate \leftarrow$AddCorrNoise$(\noisestate, \sum_{i \in \mathcal{Q}^t} \Delta^{t}_i)$}
\vspace{0.1cm}
\STATE $y^{t}, \optstate \leftarrow$ServerOpt$(y^{t-1}, \frac{1}{m}\tilde{\Delta}^t, \eta_{s}, \beta, \optstate)$ 
\ENDFOR

%\vspace{1.5ex}
\begin{mdframed}[
    linecolor=black,
    linewidth=1.5pt,
    roundcorner=4pt,
    % backgroundcolor=olive!15,
    userdefinedwidth=\linewidth,
]
\FUNCTION{ClientUpdate($i$, $x_i$)}
\STATE $\mathcal{G} \leftarrow $ (batches of user $i$'s local data)
\FOR{batch $g \in \mathcal{G}$}
\STATE $y_i \leftarrow y_i - \eta_c \nabla \ell(y_i;g)$
\ENDFOR
\STATE $\Delta \leftarrow y_i - y_i^{(0)}$
\STATE \colorbox{blue!25}{$\Delta' \leftarrow \Delta \cdot \min{\left(1, \frac{\zeta}{||\Delta||}\right)}$}
\STATE \textbf{return} $\Delta'$
\ENDFUNCTION
\end{mdframed}
\end{multicols}
\end{algorithmic}
\vspace{-0.2cm}
\end{algorithm}
}

\ifthenelse{\boolean{acl}}{\textbf{FL with DP}}{\subsection{Federated Learning with Differential Privacy (FL with DP)}}
We apply the generalized Federated Averaging (FedAvg) algorithm~\citep{mcmahan2017fedavg,wang2021fieldguide}, as shown in \cref{algo:dpfl} \ifthenelse{\boolean{acl}}{of \cref{app:background}}{}. 
FedAvg is the most common algorithm in cross-device FL systems. In a training round $t$ of total $n$ rounds, the server broadcasts a global model $y^t$ to a subset of clients; each client $i$ then updates their local model $y_i$ by SGD, and sends back the model delta; the model deltas are aggregated and used as a pseudo gradient on the server to update the global model. DP is achieved by clipping the $l_2$ norm of the model delta to control the sensitivity (contribution of each device), and then adding noise to the aggregated deltas on the server.

While our primary focus is federated learning with decentralized data in this paper, \cref{algo:dpfl} can also be applied in datacenter to achieve user-level DP~\citep{xu2022learning,chua2024mind,charles2024fine}. When using only one batch of a single sample for gradient computation in the ClientUpdate function and \treeagg for correlated noise, \cref{algo:dpfl} coincides with the DP-FTRL algorithm described in~\citep{kairouz21practical}. The DP guarantee is determined by noise calibrated to sensitivity, which depends on clip norm noise multiplier $\sigma$, the correlated noise mechanism, total number of rounds $T$, and client participation pattern (min-sep $b$). Clip norm $\zeta$ and clip norm noise multiplier $\sigma$ are used as algorithmic hyperparameters, similar to independent noise mechanism (e.g., DP-SGD/DP-FedAvg~\citep{abadi2016deep,mcmahan18learning}). However, instead of directly applying independent Gaussian noise of standard deviation $\sigma \zeta$, correlated noise are generated to privatize model updates. 

\ifthenelse{\boolean{acl}}{\textbf{MF for DP-FTRL}}{\subsection{Matrix Factorization (MF) for DP-FTRL}}
DP-FTRL~\citep{kairouz21practical} adds correlated noise to achieve strong privacy-utility trade-offs, observing that privatizing the prefix sum of model updates are essential for privatizing the training process.
The intuition of privatizing prefix sum is easier to understand when server optimizer is SGD, as the iterative process of per-round updates is equivalent to updating with prefix sum, i.e.,  
$$y^t=y^{t-1} - \eta_s \Delta^t = y^{-1}- \eta_s \sum_{j=0}^{t} \Delta^j.$$ We can write similar formulation for additional linear operation in optimization, such as momentum in SGD. In practice, it is often easier to get privatized per-round update by (conceptually) subtracting the privatized prefix sum from two consecutive rounds, and use the privatized update $\tilde{\Delta^t}$ in various server optimizers, guaranteed by the post-processing property of DP. 

We represent the model updates for $n$ rounds as a matrix $\bfX \in \R^{n \times \mdim}$, where each row is the sum of clipped updates (i.e., $\bfX_{t,\vfm} \coloneqq \sum_{i \in \mathcal{Q}^t} \Delta_i^t \in \R^{m}$ from~\cref{algo:dpfl}),
we aim to privatize $\bfA \bfX$, where $\bfA \in \R^{n \times n}$ is a lower triangular matrix of ones, i.e., $\bfA_{i,j}=1, \forall i\leq j$ and $\bfA_{i,j}=0, \forall i > j$. Given the privatized prefix sum $\widetilde{AX}$, the privatized model update is $\tilde{\Delta}^t \gets \widetilde{\bfA \bfX}_{t,\vfm} - \widetilde{\bfA \bfX}_{t-1,\vfm}$, and \cref{algo:dpfl} is privatized because of the post-processing property of DP. \citet{kairouz21practical} adopts the \treeagg\ifthenelse{\boolean{acl}}{}{\citep{dwork10continual,chan11continual,honaker15efficient}} mechanism to privatize $\bfA \bfX$. Recent work suggest a general matrix factorization framework\ifthenelse{\boolean{acl}}{~\citep{choquette2023amplified}}{~\citep{denisov22matfact, choquette2023amplified}} can be used to achieve even stronger privacy-utility trade-offs, and both \treeagg-DP-FTRL and the popular DP-SGD algorithm~\citep{abadi2016deep} are special cases in the MF-DP-FTRL framework. MF mechanism considers the factorization of $\bfA=\bfB \bfC$ and privatizes $\bfC \bfX$ by adding independent noise $\bfZ \in \R^{n \times m}$ with standard deviation $\sigma\zeta$.  
We can use the (pseudo-)inverse of the $\bfC$ matrix to generate the correlated noise in the streaming setting~\citep{choquette2023amplified},
\ifthenelse{\boolean{acl}}{
%{\small
\begin{align}
\widetilde{\bfA \bfX} & = \bfB(\bfC \bfX + \zeta \bfZ) = \bfA \bfX + \zeta \bfC^{-1} \bfZ \nonumber \\ \Rightarrow \tilde{\Delta}^{t} & = \Delta^{t} + (\zeta \bfC^{-1} \bfZ)_{t, \vfm} \label{eq:matrixmech}
\end{align}
%}%%
}{
\begin{equation}\label{eq:matrixmech}
\widetilde{\bfA \bfX} = \bfB(\bfC \bfX + \zeta \bfZ) = \bfA (\bfX + \zeta \bfC^{-1} \bfZ) \Rightarrow \tilde{\Delta}^{t} = \Delta^{t} + (\zeta \bfC^{-1} \bfZ)_{t, \vfm}
\end{equation}
}
\cref{eq:matrixmech} also suggests the alternative interpretation of correlated noise in DP-FTRL: at round t, the noise added in previous rounds can be cancelled when $\bfC^{-1}\idx{t}{\vfm}$ is negative.   

\ifthenelse{\boolean{acl}}{
\treeagg can be written in MF form, and the stronger variance reduction variant (\tafull)~\citep{honaker15efficient} is equivalent to setting $\bfB=\bfA \bfC^{-1} \in \R^{2^{l-1} \times (2^l-1)}$ by computing the Moore-Penrose pseudoinverse of $\bfC$~\citep{denisov22matfact}. However, the MF-based\tafull does not have a memory-efficient implementation, and consumes $n m$ memory. In this paper, we consider memory-efficient \treeagg~\citep{kairouz21practical} that is widely used in industry~\citep{xu2023federated}, and \tafull~\citep{denisov22matfact} that achieves better privacy-utility trade-off but less memory and computation efficient. \bandmf~\citep{choquette2023amplified} is the state-of-the-art for FL, which optimizes matrices with estimated min-sep bands. More related work  with detailed discussion are in \cref{app:ta_mf}.
}{
\paragraph{\treeagg} can be written in MF form by recursively defining $\bfC^l \in \{0 ,1\}^{(2^l-1) \times 2^{l-1}}$ $l=\ceil{\log_2 n}$, as $\bfC^1 = [1]$, $\bfC^l = [[\bfC^{l-1}, \bm{0}], [\bm{0}, \bfC^{l-1}], [\bm{1}, \bm{1}]]$, where each row $\bfC_{i,\vfm}$ represents a node in the binary tree and the ones in $\bfC_{i,\vfm}$ represent the leaf nodes for a subtree. After adding noise $\bfZ$ to every tree node, vanilla \treeagg uses matrix $\bfB$ to selects and aggregates tree nodes to privatize the prefix sum, i.e., $\bfB \in \{0, 1\}^{2^{l-1} \times (2^l-1) }$ has $\bfB_{i,j}=1, \forall j=2^{k+1}-1, k \in \kappa, i = \sum_{k \in \kappa} 2^k,$ otherwise $\bfB_{i,j}=0$. Several schemes improve vanilla binary \treeagg for prefix sums appear in the literature.  \citet{kairouz21practical} efficiently implemented \treeagg with partial variance reduction~\citep{honaker15efficient}, which leverages the recursive structure of $C$ and only needs $\ceil{\log_2 n}m$ memory to generate correlated noise. The full variance reduction trick~\citep{honaker15efficient} can further improve the performance and is equivalent to setting $\bfB=\bfA \bfC^{-1} \in \R^{2^{l-1} \times (2^l-1)}$ by computing the Moore-Penrose pseudoinverse of $\bfC$~\citep{denisov22matfact} (we use an abuse of notation $\bfC^{-1}$ for pseudoinverse). However, the full variance reduction \treeagg (\tafull) does not have a memory-efficient implementation, and consumes $n m$ memory. Another variant~\citep{andersson24smooth} is more memory-efficient, but achieves suboptimal performance compared to MF approaches~\citep{fichtenberger2022constant}. In this paper, we primarily consider \treeagg~\citep{kairouz21practical} that is widely used in industry~\citep{xu2023federated}, and \tafull~\citep{denisov22matfact} that achieves better privacy-utility trade-off but less memory and computation efficient, to represent the tree aggregation mechanisms.

\paragraph{\bandmf}
\citet{choquette2023amplified} exploits the banded structure, i.e., $\bfC \in \R^{n\times n}$ where $\bfC_{i,j}=0, \forall |i-j| \geq \hat{b}$, to simplify the optimization and privacy accounting for MF mechanisms. \bandmf successfully applied MF mechanisms to the FL system for the first time. When fixing all the other configurations for training a production LM, \bandmf improved the DP guarantee from $\rho=0.52$-zCDP by \treeagg to $\rho=0.24$-zCDP. However, \bandmf has to estimate the band size $\hat{b}$ and total rounds $n$ for optimizing matrices before training, and the performance quickly drops when the actual min-sep $b$ in FL training is smaller than $\hat{b}$, or the training round is more than $n$. \bandmf improves memory usage of MF from $n\times m$ to $\hat{b} \times m$ for correlated noise, but the typical value of min-sep $b$ in FL is still hundreds to thousands for strong DP guarantees. More recently, \bandfhu~\citep{kalinin2024banded} and \bandtoep~\citep{mckenna2024scaling} optimize banded Toeplitz matrices for larger $\nd$ and exploit Toeplitz structure for computation efficiency, but they have not been shown to outperform \bandmf in the FL setting.
}

\subsection{\blt Mechanisms in DP-FTRL} \label{sec:background_blt}
We now consider lower-triangular Toeplitz matrix in MF mechanism, i.e., $\bfC := \LtToep(c) \in \R^\dimdim$ where $\bfC\idx{i}{j}=c_{i-j}, \forall i\leq j$ otherwise $C\idx{i}{j}=0$. Buffered-linear Toeplitz (\blt) matrices \citep{mcmahan2024efficient} parameterize $\bfC$ by $\theta \in (0, 1]^\nbuf$ (the ``buffer decay'' parameters) and non-negative $\omega \in \R_+^\nbuf$ (the ``output scale'' parameters), where the Toeplitz coefficients are given by
\begin{equation} \label{eq:closedcoefs}
c_i = 
\begin{cases}
1  & i = 0 \\
\sum_{j \in [\nbuf]} \omega_j \theta_j^{i-1} & i > 0.
\end{cases}
\end{equation}

% {\color{red}
% The $\bltm(\omega, \theta)$ matrices have many useful properties,
% %\footnote{\blt matrices are closely connected to both rational functions and constant-recurrent sequences, see \citet[Sec. 3.1]{mcmahan2024efficient}. These deep mathematical roots can be highly useful, but are not necessary for our presentation.}
% most importantly for our purposes:
% \begin{enumerate*}%[label=\color{purple}(\arabic*)]
%     \item Streaming multiplication by $\bfC$, $\bfZ = \bfC \bfY$ for $\bfZ \in \R^{\nd \times \mdim}$ can be computed efficiently \citep{mcmahan2024efficient}; for completeness we give the specialization in \cref{alg:multC}.\footnote{This specialization follows by applying their Alg. 1 to the $W$ matrix of their Sec. 5.1} More relevant for our setting, a simple modification of this algorithm allow for multiplication by $\bfC\inv$, \cref{alg:multCinv}. That is, given the parameters $(\theta, \omega)$  defining $\bfC = \bltm(\theta, \omega)$, we can also efficiently compute the multiplication $\bfC\inv\bfZ$, which is precisely the key noise generation step of \cref{eq:matrixmech}.
    
%     More precisely we can compute the entries $(\bfC\inv \bfZ)\idx{i}{:}$ for $i \in [\nd]$ in streaming fashion using only $\calO(\nbuf \mdim)$ memory and $\calO(\nbuf \mdim)$ time per iterate $i$, without ever materializing $\bfC$, $\bfC\inv$, or $\bfZ$ (one need only materialize $\bfZ\idx{i}{:}$ on iteration $i$. That is, the \blt requires only $\nbuf$ memory buffers each equal in size to one row of $\bfZ$; hence, we refer to $\nbuf$ as the number of buffers, and term $\bfC$ a $\nbuf$-buffer \blt. Note \cref{alg:multCinv} did not appear in \citet{mcmahan2024efficient} and is new to this work.
    
%     \item Given a $\bfC = \bltm(\omega, \theta)$, a $\nbuf$-buffer \blt, then $\bfC\inv = \bltm(\homega, \hth)$, another $\nbuf$-buffer \blt matrix, and so in particular we can also efficiently compute its Toeplitz coefficients using \cref{eq:closedcoefs} applied to $(\homega, \hth)$.
% \end{enumerate*}

% \textbf{Remark.} The above results give two distinct methods for efficiently multiplying by $\bfC\inv$: we can apply \cref{alg:multCinv} directly to $(\theta, \omega)$, or we can derive the BLT parameterization of $\bfC\inv$, finding $(\hth, \homega)$ such that $\bfC\inv = \bltm(\hth, \homega)$, and then apply \cref{alg:multC} to $(\hth, \homega)$. We prefer to use \cref{alg:multCinv} directly, as this avoids the need to derive the parameters $(\hth, \homega)$ (which requires somewhat numerically subtle computations, e.g. polynomial root-finding or computing an SVD) in the privacy-critical code: we can easily compute sensitivity for $\bfC$ and multiply by $\bfC\inv$ using only the parameters $(\omega, \theta)$.
% }
% %%%% End Old stuff to be re-integrated


The $\bltm(\omega, \theta)$ matrices have many useful properties\ifthenelse{\boolean{acl}}{}{~\citep{mcmahan2024efficient}}, most importantly for our purposes:
\begin{enumerate*}[label=\color{purple}(\arabic*)]
    \item Streaming multiplication by $\bfC$ ($\bfZ = \bfC \hat\bfZ$ for $\bfZ,\hat\bfZ \in \R^{\nd \times \mdim}$) can be computed efficiently using only $\calO(\nbuf \mdim)$ memory and $\calO(\nbuf \mdim)$ time per round $t$, without fully materializing $\bfC$, $\bfZ$, or $\hat\bfZ$. Hence $\bfC$ is referred as a $\nbuf$-buffer \blt.
    \item The inverse of a $\nbuf$-buffer \blt ($\bfC = \bltm(\omega, \theta)$) is another $\nbuf$-buffer \blt ($\bfC\inv = \bltm(\homega, \hth)$), and we can efficiently compute Toeplitz coefficients of $\bfC\inv$ using \cref{eq:closedcoefs} applied to $(\homega, \hth)$.
\end{enumerate*}
We now derive the correlated noise generation schema for $\bfC\inv\bfZ$ in \cref{eq:matrixmech} based on the BLT properties. We can first derive the BLT parameters $(\hth, \homega)$ of $\bfC\inv$ such that $\bfC\inv = \bltm(\hth, \homega)$, and then generate the correlated noise based on $(\hth, \homega)$ in streaming setting. However, we show a simpler alternative that directly uses the BLT parameters $(\theta, \omega)$ to generate streaming correlated noise $\hat\bfZ$. 

Applying the parameterization in \cref{eq:closedcoefs} \ifthenelse{\boolean{acl}}{}{(also appear in \citep[Sec 5]{mcmahan2024efficient})} to \citep[Alg 1]{mcmahan2024efficient}, and initializing buffers $\bfS_{-1} \gets \mathbf{0} \in \R^{\nbuf \times \mdim}$, we can efficiently compute $\bfZ_{t,\vfm}$ from $\hat\bfZ_{t,\vfm}$ in the streaming setting,
  $
  \bfZ_{t,\vfm} = \hat\bfZ_{t,\vfm} + \outp^T \bfS_{t-1}, 
  \bfS_{t} = \diag(\theta) \bfS_{t-1}  + \bm{1}_\nbuf \hat\bfZ_{t,\vfm}
  $. We rearrange the update equations to get $\hat\bfZ$ from $\bfZ$ and $\bfS$, 
  \ifthenelse{\boolean{acl}}{
  %{\small
  \begin{align}
  \hat\bfZ_{t,\vfm}  = \bfZ_{t,\vfm} - \outp^T \bfS_{t-1}, \nonumber \\ 
  \bfS_{t} = \diag(\theta) \bfS_{t-1}  + \bm{1}_\nbuf \hat\bfZ_{t,\vfm}. \label{eq:blt_noise}
  \end{align}
  %}%%
  }{
  \begin{align}
  \hat\bfZ_{t,\vfm}  = \bfZ_{t,\vfm} - \outp^T \bfS_{t-1},  
  \bfS_{t} = \diag(\theta) \bfS_{t-1}  + \bm{1}_\nbuf \hat\bfZ_{t,\vfm}. \label{eq:blt_noise}
  \end{align}
  }
  To efficiently generate correlated noise $\hat\bfZ_{t,\vfm}$ at round $t$, we only need to materialize the independent noise $\bfZ_{t,\vfm} \in \R^{1 \times m}$ and use the buffers $\bfS_{t-1} \in \R^{\nbuf \times \mdim}$ in \cref{eq:blt_noise}. The efficient correlated noise generation parameterized by BLT parameters $(\theta, \omega)$ for $\bfC\inv$ (instead of $\bfC$) did not appear in \citet{mcmahan2024efficient} and is new to this work. \cref{eq:blt_noise} intuitively shows the noise cancellation view of correlated noise, where the previous noises are tracked in the states decaying with $\theta \in (0, 1]$, and then subtracted in the current round after scaling with $\omega$. 
  
  For completeness, we provide the streaming multiplication algorithm of $\bfZ = \bfC \hat \bfZ$ and $\hat \bfZ = C^{-1} \bfZ$ for $\bfC = \bltm(\theta, \omega)$ in \cref{alg:multC} and \cref{alg:multCinv}, respectively. 
  \cref{alg:multC} is a direct application of \citep[Alg 1]{mcmahan2024efficient} and only used to derive \cref{alg:multCinv}, which is our streaming algorithm for generating correlated noise with \blts using only $dm$ memory. 
  Finally, we apply the streaming correlated noise $\hat\bfZ$ by BLT from \cref{alg:multCinv} (using \cref{eq:blt_noise})  in \cref{algo:dpfl} for the BLT-DP-FTRL algorithm, i.e., 
  \begin{equation}
      \tilde{\Delta}^t \leftarrow \sum_{i \in \mathcal{Q}^t} \Delta^{t}_i + \hat\bfZ_{t,\vfm}.
  \end{equation}

 
\ifthenelse{\boolean{acl}}{}{ \begin{minipage}[t]{0.49\textwidth}
\begin{algorithm}[H]
\caption{\small Stream Mult. by $\bltm(\theta, \omega)$~\citep{mcmahan2024efficient}} \label{alg:multC}
\begin{algorithmic}%[1]
\STATE \textbf{Input:} 
\STATE \myind Input stream $\hat \bfZ \in \R^{\nd \times \mdim}$
\STATE \myind $\theta \in \R^\nbuf, \omega \in \R^\nbuf$ for $\bfC = \bltm(\theta, \omega)$
\STATE \textbf{Output:}
\STATE \myind The rows $\bfZ\idx{t}{:}$ of $\bfZ = \bfC \hat \bfZ$
\STATE
\STATE Initialize buffers $\Buf_{-1} \assign \mathbf{0} \in \R^{\nbuf \times \mdim}$
\FOR{$t=0, \dots, n - 1$}
  \STATE  $\bfZ_{t,\vfm} = \hat\bfZ_{t,\vfm} + \outp^T \bfS_{t-1}$ 
  \STATEComment Decay each buffer by $\theta$ and add $\hat\bfZ_{t,\vfm}$ to each
  \STATE $\bfS_{t} = \diag(\theta) \bfS_{t-1}  + \bm{1}_\nbuf \hat\bfZ_{t,\vfm}$
  \STATE Output $\bfZ_{t,\vfm}$
\ENDFOR
\end{algorithmic}
\end{algorithm}
\end{minipage}
%
\hfill
%
\begin{minipage}[t]{0.49\textwidth}
\begin{algorithm}[H]
\caption{\small Stream Mult. by $\bltm^{-1}(\theta, \omega)$} \label{alg:multCinv}
\begin{algorithmic}%[1]
\STATE \textbf{Input:} 
\STATE \myind Input stream $\bfZ \in \R^{\nd \times \mdim}$
\STATE \myind $\theta \in \R^\nbuf, \omega \in \R^\nbuf$ for $\bfC = \bltm(\theta, \omega)$
\STATE \textbf{Output:}
\STATE \myind The rows $\hat \bfZ\idx{t}{:}$ of $\hat \bfZ = \bfC^{-1} \bfZ$

\STATE
\STATE Initialize buffers $\Buf_{-1} \assign \mathbf{0} \in \R^{\nbuf \times \mdim}$
\FOR{$t=0, \dots, n - 1$}
  \STATE $\hat\bfZ_{t,\vfm}  = \bfZ_{t,\vfm} - \outp^T \bfS_{t-1}$
  \STATEComment The buffer update is the same as \cref{alg:multC}
  \STATE $  \bfS_{t} = \diag(\theta) \bfS_{t-1}  + \bm{1}_\nbuf \hat\bfZ_{t,\vfm}$
  \STATE Output $\hat \bfZ\idx{t}{:}$
\ENDFOR
\end{algorithmic}
\end{algorithm}
\end{minipage}}

\ifthenelse{\boolean{acl}}{}{\begin{table}[t]
\renewcommand{\arraystretch}{1.4}
\centering
\small
\begin{tabular}{p{0.9in} p{0.7in} p{0.8in} p{0.9in} p{0.5in} p{0.4in}}
\toprule
Mech & Loss & Mech. opt. \newline cost & Memory \newline overhead  & Noise Added & $(n,\minsep)$-fragility \\
\midrule
\blt (ours)
  & \maxloss/ \rmsloss
  & Low \newline ($\calO(n)$ )  
  & Low \newline ($\sim 4 \times \mdim$  )         
  & Low
  & Low \\
\bandmf\ifthenelse{\boolean{acl}}{}{\citep{choquette2023amplified}}
  & \rmsloss
  & High \newline ($\calO(n^2)$)
  & High \newline ($\bands \times \mdim$ )     
  & Lowest 
  & Med \\
\bandtoep\ifthenelse{\boolean{acl}}{}{\citep{mckenna2024scaling}}
  & \rmsloss
  & Low  \newline ($\calO(n)$)
  & High \newline ($\bands \times \mdim$)
  & Low 
  & Med \\
% TODO - double check on memory overhead for Full Honaker.
\treeagg\ifthenelse{\boolean{acl}}{}{\citep{kairouz21practical}}  
  & \rmsloss\!\!*
  & Lowest \newline (predefined)
  & Med \newline  ($\ceil{\log_2(\nd)} \times \mdim$) 
  & High 
  & Low  \\
\bottomrule
\end{tabular}
\caption{\small Summary of mechanisms considered, evaluated in terms of: \emph{mechanism optimization cost}, how expensive is it to compute the mechanism $\bfC$; $\calO(\cdot)$ gives the cost of a single gradient calculation. The next two columns relate to the deployment of the mechanism in a ML training system: \emph{memory overhead} is the additional state (as a multiple of the model dimension $\mdim$) that the server needs to maintain; the per-round runtime cost is also proportional to this value. The \emph{noise added} is categorized subjectively, for details see examples in ~\cref{fig:rmse}. $(n, \minsep)$-fragility reflects the degree to which the mechanisms performance degrades when total number of rounds $n$ and min-sep $\minsep$ are poorly estimated, see discussion on quantitative results in \cref{fig:rmse_k3,fig:rmse_nkb,fig:rmse_vary_n,fig:other_mechs}.  The conclusion:  \textbf{\blts perform well on all aspects.}
} \label{tab:mechanism_summary}
\end{table}


\subsection{Comparing DP-FTRL Mechanisms}
We summarize DP-FTRL mechanisms in \cref{tab:mechanism_summary} \ifthenelse{\boolean{acl}}{in \cref{app:blt-dp-ftrl}}{} and show the advantages of BLT-DP-FTRL. Our BLT mechanism can optimize either MaxLoss or RmsLoss for generating correlated noise (detailed in \cref{sec:opt_blt} following~\citep{mcmahan2024efficient}), while previous MF mechanisms in practice primarily consider RmsLoss~\citep{choquette2023amplified,mckenna2024scaling}. It is possible to extend the previous mechanisms to use MaxLoss, while it is still an open problem which loss format corresponds better with learning performance when running with the DP-FTRL algorithms. \treeagg, especially \tafull, is equivalent to considering RmsLoss though the mechanism is predefined without explicit optimization cost; and we present the lower memory overhead ($\ceil{\log_2(\nd)} \times \mdim$) for \treeagg while \tafull without an efficient algorithm yet actually needs . 

BLTs achieve better privacy-utility trade-offs than \tafull in simulation benchmark experiments (see \cref{sec:exp}), and clearly outperforms \treeagg in production cross-device FL experiments (see \cref{sec:gboard}), as lower noise is added in BLTs. While \bandmf\citep{choquette2023amplified} can add lowest noise (measured by RmsLoss in \cref{fig:rmse}), BLTs have lower mechanism optimization cost and memory overhead. Moreover, \cref{sec:exp,sec:gboard} show the learning performance of BLTs are often comparable with \bandmf under the same privacy guarantee in practical settings, though BLTs' RmsLoss is slightly worse. The memory overhead of BLTs is $d\times m$ where we empirically observe that buffer size $d=4$ achieves low losses and further increasing $d$ does not further improve in our empirical settings of total rounds $n$ and min-sep $b$. The BLT memory overhead of $d \sim 4$ is smaller than \treeagg where $\ceil{\log_2(n)} \sim 11$, and much smaller than typical $\hat b \sim X00$ for \bandmf and \bandtoep. \bandtoep~\citep{mckenna2024scaling} suggested small $\hat b$ is preferred when using amplification by sampling in the many participation settings; however, sampling is generally not possible in practical cross-device FL systems. 

As shown in \cref{fig:rmse_k3,fig:rmse_nkb,fig:rmse_vary_n,fig:other_mechs}, BLT is also more robust than banded MF mechanisms when number of total rounds $n$ and min-sep $b$ are not accurately estimated. Specifically, it is unclear how to run Banded MF mechanisms beyond the estimated $n$ after optimizing the $\bfC \in \R^{n, n}$ matrix for correlated noise. Optimizing $\bfC \in \R^{n, n}$ for a much larger $n$ and truncating it to the actual number of training rounds can achieve good privacy-utility trade-offs, but encounter non-trivial mechanism optimization cost. BandMF performance degrades fast when the actual min-sep $b$ is smaller than the estimated band $\hat b$, but the stronger DP guarantees are generally achieved when $\hat b$ is large. Hence the tricky task of estimating min-sep $b$ is more important for \bandmf. In general, BLT-DP-FTRL is competitive for achieving state-of-the-art privacy-utility trade-off (compared to \bandmf), while maintains ease to use in practice (compared to \treeagg).  }
